\chapter{Anomaly Detection}
\label{ch:anomaly}

\chapterguide%
  {\chapref{ch:unsupervised}, distributions and quantiles from \chapref{ch:foundations}, plus clustering and density ideas from \chapref{ch:clustering}.}
  {flag unusual observations with statistical rules, local density, and isolation forests; distinguish outlier from novelty detection; choose operational thresholds; and evaluate ranked anomalies responsibly.}

This chapter completes the core unsupervised sequence with \term{anomaly detection}: ranking observations that do not fit the reference pattern --- fraudulent transactions, failing sensors, intrusions, or data-entry errors.

\begin{firstread}
An anomaly detector is an alarm system, not a judge. It ranks cases by how unusual they appear so a person or later process can investigate them.
\end{firstread}

\section{Anomaly detection}

Two closely related setups share the name:

\begin{mltable}
\begin{tabularx}{\linewidth}{@{}>{\bfseries\leavevmode\color{MLNavy}}p{0.22\linewidth}p{0.36\linewidth}X@{}}
\toprule
\rowcolor{MLPaleBlue!92}
Setup & Training data & Goal \\
\midrule
Outlier detection & The training data itself is contaminated with anomalies. & Find the anomalous observations already present in the data. \\
Novelty detection & The training data contains clean ``normal'' behavior. & Flag future observations that deviate from the learned normal pattern. \\
\bottomrule
\end{tabularx}
\end{mltable}

Labels are scarce almost by definition --- anomalies are rare, expensive to confirm, and unlike one another --- so most methods learn what \emph{normal} looks like and score each observation by how far it strays.

\subsection{Statistical rules: z-scores and IQR fences}

For a single roughly bell-shaped feature, the \term{z-score} measures how many standard deviations an observation sits from the mean:
\[
  z=\frac{x-\bar{x}}{s},
\]
with $|z|>3$ a customary alarm threshold --- under a normal distribution, only about $0.3\%$ of values land that far out.

\begin{workedexample}
A temperature sensor averages $\bar{x}=50$ with standard deviation $s=5$. A reading of $68$ has $z=(68-50)/5=3.6$: flagged. A reading of $57$ has $z=1.4$: unremarkable.
\end{workedexample}

The z-score has an ironic weakness: the anomalies it hunts inflate the very mean and standard deviation it depends on, hiding themselves. Robust variants substitute the median for the mean and a scaled median absolute deviation for $s$. The same robustness motivates the \term{IQR fences} from \chapref{ch:statistics}: with quartiles $Q_1$ and $Q_3$ and $\operatorname{IQR}=Q_3-Q_1$, flag anything outside
\[
  [\,Q_1-1.5\operatorname{IQR},\ Q_3+1.5\operatorname{IQR}\,],
\]
the rule behind the whiskers of a box plot. If delivery times have $Q_1=20$ and $Q_3=30$ minutes, the fences sit at $5$ and $45$ minutes, and a 60-minute delivery is flagged --- no distributional assumptions required.

Per-feature rules are the right first tool, but they miss \emph{combinations}: a heart rate of 150 is normal, ``asleep'' is normal, and the two together are an emergency. Multivariate methods exist for exactly this.

\subsection{Distance and local density}

A direct multivariate score is the distance to an observation's $k$-th nearest neighbor: normal points sit in crowds, anomalies sit alone. Its refinement, the \term{local outlier factor (LOF)}, fixes a subtle failure of global distance. In a dataset containing one tight cluster and one naturally diffuse cluster, every member of the diffuse cluster looks ``far from neighbors'' by a global yardstick. LOF instead compares each point's local density with the densities around its own neighbors: a score near 1 means ``about as crowded as my neighborhood expects,'' while scores well above 1 mean ``markedly emptier around me than around my neighbors'' --- an outlier \emph{relative to its local context}.

\subsection{Model-based scores}

Any density estimate can support anomaly scoring by thresholding likelihood: fit a Gaussian or a Gaussian mixture (\chapref{ch:clustering}) to reference data and flag observations the model considers improbable. The \term{one-class SVM} adapts \chapref{ch:svm}'s machinery to the same end, learning a frontier around the bulk of the data and flagging observations outside it.

\subsection{Isolation forests}
\label{sec:isolation-forest}

The isolation forest inverts the usual logic: instead of modeling normal points, it measures how easy each point is to \emph{isolate}. Build many random trees, each splitting on a random feature at a random threshold between the observed minimum and maximum, and record the depth at which each observation ends up alone.

\begin{intuition}
Play twenty questions against the dataset. A point deep inside the crowd needs many random questions before it stands alone --- every split keeps plenty of company. A point far outside needs only a couple: ``is $x_1>90$?'' and it is already by itself. Average that question-count over many random trees; a short normalized path length drives a more anomalous score.
\end{intuition}

Isolation forests need no distance function and can be an efficient general-purpose baseline. Their \term{contamination} setting calibrates the score threshold to an assumed or operational anomaly fraction; it does not discover the true prevalence.

\subsection{Evaluating and using anomaly detectors}

\begin{itemize}
  \item With few or no labels, evaluation leans on ranked review: send the top-scored cases to a human and measure precision among the top $k$; when some labels exist, precision-recall analysis (\chapref{ch:evaluation}) is usually more useful than accuracy under extreme imbalance.
  \item For distance- and kernel-based detectors, apply the scaling rule from \secref{sec:scaling}.
  \item Choose thresholds from consequences --- how many alarms per day can be investigated? --- not from statistical custom.
  \item An anomaly is not automatically an error. A rare observation may be invalid, or it may be the most important case in the dataset: the fraud, the failing part, the discovery. Investigate before deleting.
\end{itemize}

\section{The shape of an anomaly problem}

Anomaly detection differs from classification in one important way: you usually cannot see many examples of the thing you are looking for.

\begin{mlfigure}
\begin{tikzpicture}[font=\scriptsize]
\begin{scope}
\draw[draw=MLMuted,fill=white] (0,0) rectangle (4.2,2.9);
\foreach \p in {(1.5,1.4),(1.8,1.6),(1.6,1.1),(2.0,1.3),(1.3,1.7),(2.2,1.6),(1.9,0.9),
                (1.4,1.25),(2.1,1.1),(1.7,1.45),(2.3,1.35),(1.55,0.85),(2.05,1.65)}
  \fill[MLBlue] \p circle (2.4pt);
\draw[MLTeal,thick,dashed] (1.8,1.32) circle (0.95);
\fill[MLRed] (3.5,2.4) circle (2.8pt);
\fill[MLRed] (0.5,0.4) circle (2.8pt);
\node[color=MLRed,anchor=west] at (3.0,2.72) {anomalies};
\node[color=MLTeal] at (1.8,-0.35) {the "normal" region};
\node[font=\small\bfseries\color{MLNavy}] at (2.1,3.25) {Distance / density view};
\end{scope}
\begin{scope}[xshift=5.6cm]
\draw[draw=MLMuted,fill=white] (0,0) rectangle (4.6,2.9);
\draw[-{Stealth[length=2mm]},MLMuted] (0.3,0.45) -- (4.35,0.45);
\node[color=MLMuted,anchor=north] at (2.3,0.35) {anomaly score};
\draw[MLBlue,very thick] plot[smooth,domain=0.4:4.3,samples=60]
  (\x, {0.45+2.0*exp(-((\x-1.3)^2)/0.28)});
\draw[MLRed,very thick,dashed] (3.05,0.45) -- (3.05,2.6);
\node[color=MLRed,anchor=west,align=left] at (3.15,2.2) {threshold\\you choose};
\node[color=MLGreen,anchor=east] at (2.2,1.05) {normal};
\node[color=MLRed,anchor=west] at (3.3,0.85) {flagged};
\node[font=\small\bfseries\color{MLNavy}] at (2.3,3.25) {Score-and-threshold view};
\end{scope}
\end{tikzpicture}
\end{mlfigure}

\begin{keyidea}
Every anomaly detector produces a \emph{score}, not a verdict. Turning that score into a decision requires a threshold, and the threshold is a business choice: how many false alarms can your team investigate per day? A detector that flags 10\,000 transactions is useless if only 50 can be reviewed.
\end{keyidea}

\begin{analogy}
Classification is being shown a thousand photographs of cats and a thousand of dogs and learning to tell them apart. Anomaly detection is being shown a thousand ordinary days at a shop and asked to notice when something is off -- without ever having been shown a burglary. You learn what normal looks like, and treat anything unlike it as suspicious.
\end{analogy}


\section{From anomaly score to investigation}

The technical model is only half of an anomaly system. The other half is the review process that consumes its ranked output.

\begin{mltable}
\begin{tabularx}{\linewidth}{@{}>{\bfseries\leavevmode\color{MLNavy}}p{0.35\linewidth}X@{}}
\toprule
\rowcolor{MLPaleBlue!92}
Question & Practical consequence \\
\midrule
How many cases can be reviewed? & Choose a top-$k$ or score threshold that respects review capacity. \\
What happens after a flag? & Define whether the case is blocked, queued, inspected, or merely logged. \\
Which false alert is costly? & Track precision among reviewed cases and the burden on users or analysts. \\
Can normal behavior change? & Monitor score distributions and retrain or recalibrate when the reference population drifts. \\
Can labels arrive later? & Use confirmed outcomes to evaluate ranking quality and eventually consider a supervised model. \\
\bottomrule
\end{tabularx}
\end{mltable}

\begin{modelcard}{Anomaly Detection - Quick Guide}
\begin{tabularx}{\linewidth}{@{}>{\bfseries\leavevmode\color{MLNavy}}p{0.25\linewidth}X@{}}
Anomaly detection asks & Which observations are unusually inconsistent with a reference pattern? \\
Useful methods & Robust statistical rules, distance/local-density methods, isolation forest, and one-class models. \\
Output & An anomaly score or ranking; a later policy converts that score into an alert. \\
Evaluation & Precision among reviewed cases, PR AUC when labels exist, stability, and downstream review value. \\
Main risks & Rare does not always mean harmful; distribution shift can make yesterday's normality model obsolete. \\
\end{tabularx}
\end{modelcard}

% ==== EXPANDED EDITION ADDITIONS ====

\section{Kaggle implementation: anomaly ranking on small fraud data}
\label{sec:kaggle-credit-anomaly}

\begin{kaggleproject}{Credit-card transactions --- Isolation Forest without label-aware fitting}
\kagglesource{https://www.kaggle.com/datasets/miadul/credit-card-fraud-detection-dataset}{Credit Card Fraud Detection Dataset}
\textbf{Question.} Can an unsupervised anomaly score rank rare fraudulent transactions above ordinary transactions even though \texttt{is\_fraud} is not supplied to \texttt{fit}?
\end{kaggleproject}

This public dataset contains 10,000 transactions and is only a few hundred kilobytes, making it much easier to keep beside the course notebooks than the commonly used 150 MB credit-card dataset. The Isolation Forest fits only on transaction-behavior features; the fraud label is held back from training and used afterward to evaluate the anomaly ranking with ROC AUC, average precision, and a concrete top-$k$ review budget. Average precision is also compared with fraud prevalence so the ranking gain is interpreted relative to the no-skill precision level.

\textbf{Before running.} Predict the qualitative pattern you expect from the experiment before you look at the output or plot.

\begin{labmode}
Stage 4B: decide the review budget and ranking metrics before fitting. The held-back fraud labels are for evaluation only, not for fitting the anomaly detector.
\end{labmode}

\lstinputlisting[style=mlpython]{all_experiments/lab17-credit_card_anomaly.py}


\begin{pitfall}
An anomaly detector finds observations that are unusual under its representation; unusual does not automatically mean fraudulent, defective, malicious, or important. The semantic meaning comes from external validation.
\end{pitfall}

\begin{labdeliverable}
Submit fraud prevalence, ROC AUC, average precision, the ratio of average precision to prevalence, and precision among the top-ranked review budget you choose. Include one ranked-case table and explain why anomaly score, fraud probability, and semantic proof of fraud are different claims.
\end{labdeliverable}



\begin{quickcheck}
\begin{enumerate}
  \item Why is an anomaly score not automatically a probability of fraud, failure, or error?
  \item How does a fixed review capacity turn anomaly detection into a ranking problem?
  \item What is the difference between detecting outliers in the training sample and novelty detection for future observations?
\end{enumerate}
\end{quickcheck}

\begin{chapterrecap}
\begin{itemize}
  \item Anomaly detection assigns an unusualness score rather than discovering a guaranteed error or fraud case.
  \item Robust rules use medians, IQRs, distances, local density, or isolation behavior.
  \item Evaluation requires labeled cases when available and careful review of false alerts.
\end{itemize}
\end{chapterrecap}
